\chapter{ΕΙΣΑΓΩΓΗ}

\section{Γενικά}

Τα σύγχρονα Συστήματα Ηλεκτρικής Ενέργειας (ΣΗΕ) διέρχονται μια περίοδο θεμελιωδών μετασχηματισμών, η οποία χαρακτηρίζεται από τη σταδιακή απόσυρση των συμβατικών θερμικών μονάδων παραγωγής και τη μαζική διείσδυση σταθμών Ανανεώσιμων Πηγών Ενέργειας (ΑΠΕ). Η αντικατάσταση των κλασικών σύγχρονων γεννητριών από μονάδες παραγωγής βασισμένες σε ηλεκτρονικά ισχύος (όπως οι ανεμογεννήτριες επαγωγικής γεννήτριας διπλής τροφοδότησης (DFIG) και οι μετατροπείς VSC) μεταβάλλει ριζικά τη δυναμική συμπεριφορά και την ευστάθεια του δικτύου.

Ένα σημαντικό χαρακτηριστικό της μετάβασης είναι η μείωση της συνολικής φυσικής (περιστροφικής) αδράνειας του συστήματος. Η μειωμένη αδράνεια καθιστά το ηλεκτρικό δίκτυο πιο ευάλωτο σε διαταραχές, με αποτέλεσμα την εμφάνιση εντονότερων και συχνότερων ηλεκτρομηχανικών ταλαντώσεων. Οι ταλαντώσεις αυτές διακρίνονται κυρίως σε:

\begin{itemize}
    \item Intra-area: Ταλαντώσεις τοπικού χαρακτήρα με συχνότητες 0.8Hz \textendash{} 2 Hz, οι οποίες αφορούν τη σχετική κίνηση μιας γεννήτριας ως προς το υπόλοιπο τοπικό δίκτυο.

    \item Inter-area: Ταλαντώσεις ευρείας κλίμακας με συχνότητες 0.1Hz \textendash{} 0.8Hz, οι οποίες εκδηλώνονται μεταξύ απομακρυσμένων ομάδων γεννητριών μέσω γραμμών διασύνδεσης μεγάλου μήκους.
\end{itemize}

Η ανεπαρκής απόσβεση των ταλαντώσεων αυτών ενέχει σοβαρούς κινδύνους για τη δυναμική ευστάθεια του συστήματος, καθώς μπορεί να οδηγήσει σε αστάθεια, ανεπιθύμητες ενεργοποιήσεις συστημάτων προστασίας ή ακόμη και σε εκτεταμένες διακοπές ρεύματος.

\section{Τεχνολογική Αιχμή Θέματος}

Για την αντιμετώπιση των προαναφερθέντων προκλήσεων, τα σύγχρονα δίκτυα αξιοποιούν Συστήματα Παρακολούθησης Ευρείας Περιοχής (WAMS) με χρήση Μονάδων Μέτρησης Φασιθετών (PMUs). Οι PMUs παρέχουν χρονικά συγχρονισμένες μετρήσεις υψηλού ρυθμού δειγματοληψίας, επιτρέποντας την παρακολούθηση της δυναμικής του συστήματος σε πραγματικό χρόνο.

Στην τεχνολογική αιχμή της ανάλυσης σημάτων WAMS βρίσκονται οι προηγμένες τεχνικές ταυτοποίησης συστημάτων και οι αλγόριθμοι μηχανικής μάθησης:

\begin{itemize}
    \item Ταυτοποίηση Ringdown \& Ambient: Η επεξεργασία σημάτων διακρίνεται σε αναλύσεις Ringdown μετά από μεγάλες διαταραχές (μέσω αλγορίθμων όπως Matrix Pencil, Prony, ERA, N4SID, Vector Fitting) και αναλύσεις Ambient κατά τη φυσιολογική λειτουργία υπό την επίδραση τυχαίων διακυμάνσεων (μέσω μεθόδων αναγνώρισης υποχώρου όπως η N4SID).

    \item Αλγόριθμοι Συσταδοποίησης: Λόγω του μεγάλου όγκου δεδομένων από πολλαπλές PMUs και της παρουσίας θορύβου ή αριθμητικών/ψευδών ιδιοτιμών, η εφαρμογή αλγορίθμων συσταδοποίησης (όπως k-Means, k-Medoids, DBSCAN, OPTICS, HDBSCAN) αποτελεί απαραίτητο εργαλείο για την αυτόματη ομαδοποίηση και την απομόνωση των πραγματικών ρυθμών του συστήματος.

    \item Ποσοτική Αξιολόγηση: Η χρήση εξειδικευμένων δεικτών, όπως ο Median Absolute Deviation (MAD) για την ακρίβεια θέσης και οι Adjusted Rand Index (ARI) και V-measure για τη συμφωνία των συστάδων, επιτρέπει τη συστηματική αξιολόγηση της αξιοπιστίας των αποτελεσμάτων.
\end{itemize}

\section{Αντικείμενο Διπλωματικής}

Το αντικείμενο της παρούσας διπλωματικής εργασίας είναι η εκτίμηση, ταυτοποίηση και συσταδοποίηση των ηλεκτρομηχανικών ταλαντώσεων σε Συστήματα Ηλεκτρικής Ενέργειας. Στο πλαίσιο αυτό, η έρευνα επικεντρώνεται αρχικά στη διεξαγωγή δυναμικών RMS προσομοιώσεων σε πρότυπα συστήματα δοκιμών (Kundur Two-Area System και IEEE 39-Bus System) για την παραγωγή και συλλογή αντιπροσωπευτικών σημάτων Ringdown και Ambient.

Για την επεξεργασία των σημάτων αυτών, εφαρμόζονται προηγμένες τεχνικές ταυτοποίησης συστημάτων. Συγκεκριμένα, αξιοποιείται η μέθοδος Matrix Pencil (MP) σε σήματα Ringdown με σκοπό την ανακατασκευή των κυματομορφών και την εκτίμηση των παραμέτρων ταλάντωσης, ενώ η μέθοδος N4SID εφαρμόζεται σε Ambient δεδομένα για την εκτίμηση των ιδιοτιμών υπό συνθήκες φυσιολογικής λειτουργίας του δικτύου. Στη συνέχεια, πραγματοποιείται η εφαρμογή και η συγκριτική αξιολόγηση επτά αλγορίθμων συσταδοποίησης, αποσκοπώντας στη συστηματική ομαδοποίηση των εκτιμώμενων ιδιοτιμών και την αποτελεσματική απόρριψη των ψευδών ρυθμών ταλάντωσης.

Η εργασία εστιάζει, παράλληλα, στην εξέταση της επίδρασης που έχει η διείσδυση αιολικών μονάδων τεχνολογίας DFIG στη μεταβολή των ηλεκτρομηχανικών ρυθμών του συστήματος, αξιολογώντας τη δυναμική συμπεριφορά του δικτύου υπό νέες συνθήκες λειτουργίας.

\section{Συμβολή}

Η κύρια συμβολή της παρούσας διπλωματικής εργασίας έγκειται στην ανάπτυξη ενός ολοκληρωμένου και συστηματικού πλαισίου ταυτοποίησης και συσταδοποίησης για τη διαχείριση δεδομένων WAMS. Η προτεινόμενη μεθοδολογία συνδυάζει προηγμένες τεχνικές ταυτοποίησης στο πεδίο του χρόνου με επτά διαφορετικούς αλγόριθμους συσταδοποίησης, προσφέροντας μια ολιστική προσέγγιση στην ανάλυση των ηλεκτρομηχανικών ταλαντώσεων.

Παράλληλα, παρέχεται μια εκτενής συγκριτική αξιολόγηση αλγορίθμων συσταδοποίησης (συγκεκριμένα των k-Means, k-Medoids, DBSCAN, OPTICS, HDBSCAN, GMM και Agglomerative Hierarchical Clustering) ως προς τη συνοχή, την κάλυψη, τη στιβαρότητα απέναντι στον θόρυβο και την ακρίβεια χαρτογράφησης των ιδιοτιμών. Για την αντικειμενική μέτρηση της απόδοσης των αλγορίθμων αυτών σε σχέση με τις θεωρητικές τιμές αναφοράς, εισάγεται η συστηματική χρήση εξειδικευμένων δεικτών αξιολόγησης, όπως ο Median Absolute Deviation (MAD), ο Adjusted Rand Index (ARI) και το V-measure.

Η εργασία συμβάλλει στη μελέτη της επίπτωσης της αυξημένης διείσδυσης ΑΠΕ στη δυναμική του δικτύου, αποτυπώνοντας ποσοτικά τη μεταβολή των ηλεκτρομηχανικών ρυθμών λόγω της αντικατάστασης συμβατικών γεννητριών από αιολικά πάρκα τεχνολογίας DFIG. Επιπρόσθετα, αξιολογείται η προσαρμοστικότητα και η αξιοπιστία των μεθόδων συσταδοποίησης υπό τις μεταβαλλόμενες συνθήκες λειτουργίας του συστήματος.

\section{Δομή της Διπλωματικής}

Η εργασία διαρθρώνεται σε έξι κεφάλαια ως εξής:

\begin{itemize}
    \item Κεφάλαιο 1 - Εισαγωγή: Παρουσιάζεται το πλαίσιο, το κίνητρο, η τεχνολογική αιχμή, το αντικείμενο, η συμβολή και η δομή της διπλωματικής εργασίας.

    \item Κεφάλαιο 2 - Θεωρητικό Υπόβαθρο: Αναλύονται οι ηλεκτρομηχανικές ταλαντώσεις, οι τεχνολογίες WAMS/PMU, οι μέθοδοι ταυτοποίησης (Matrix Pencil, N4SID κ.ά.), οι επτά αλγόριθμοι συσταδοποίησης, οι δείκτες αξιολόγησης (MAD, ARI, V-measure, BIC), καθώς και τα δυναμικά μοντέλα και τα πρότυπα συστήματα δοκιμών (Kundur Two-Area, IEEE 39-Bus).

    \item Κεφάλαιο 3 - Ανάλυση Ringdown και Ταυτοποίηση με Matrix Pencil: Περιγράφεται η παραγωγή δεδομένων Ringdown στο PowerFactory, η προεπεξεργασία των σημάτων και η εφαρμογή της μεθόδου Matrix Pencil για την ανακατασκευή κυματομορφών και την εκτίμηση ιδιοτιμών.

    \item Κεφάλαιο 4 - Ambient Ανάλυση στο Σύστημα IEEE 39-Bus: Πραγματοποιείται εκτίμηση ιδιοτιμών με τη μέθοδο N4SID σε Ambient δεδομένα και εφαρμόζονται οι επτά αλγόριθμοι συσταδοποίησης για τη συγκριτική τους αξιολόγηση.

    \item Κεφάλαιο 5 - Ambient Ανάλυση με Μονάδες DFIG: Μελετάται το τροποποιημένο σύστημα IEEE 39-Bus με διείσδυση αιολικών μονάδων DFIG, εξετάζεται η μεταβολή των ηλεκτρομηχανικών ρυθμών και αξιολογείται η απόδοση των μεθόδων συσταδοποίησης.

    \item Κεφάλαιο 6 - Συμπεράσματα: Συνοψίζονται τα κύρια συμπεράσματα της εργασίας και προτείνονται κατευθύνσεις για μελλοντική έρευνα.
\end{itemize}
